All model comparison metrics were computed on the original response scale using
the fitted mean trajectory for each model. For models that produce one curve per
trial, write this fitted value as $\widehat y_i^{[j]}$. For the single-curve
Neural ODE, $\widehat y_i^{[j]}=\widehat y_i$ for every trial $j$. Let
$$
  \mathrm{SSE}
  =
  \sum_{(i,j)\in\Omega}(y_i^{[j]}-\widehat y_i^{[j]})^2,
  \qquad
  \mathrm{SST}
  =
  \sum_{(i,j)\in\Omega}(y_i^{[j]}-\bar y)^2.
$$
The root mean squared error is
$$
    \mathrm{RMSE}
    =
    \sqrt{\frac{\mathrm{SSE}}{n}},
$$
and the coefficient of determination is
$$
    R^2
    =
    1-\frac{\mathrm{SSE}}{\mathrm{SST}}.
$$

The numerical table below reports in-sample residual diagnostics: the models
are trained on all observed entries in $\Omega$, and the same entries are used
to compute RMSE and $R^2$. This is appropriate for checking interpolation
quality, but it is not an out-of-sample generalization error. To estimate a
test MSE for this time-series problem, the split should be made by time rather
than by randomly sampling individual observations. Let
$$
  I_{\mathrm{train}}\subset \{0,\ldots,T-1\},
  \qquad
  I_{\mathrm{test}}\subset \{0,\ldots,T-1\},
  \qquad
  I_{\mathrm{train}}\cap I_{\mathrm{test}}=\varnothing,
$$
and define
$$
  \Omega_{\mathrm{train}}
  =
  \{(i,j)\in\Omega:i\in I_{\mathrm{train}}\},
  \qquad
  \Omega_{\mathrm{test}}
  =
  \{(i,j)\in\Omega:i\in I_{\mathrm{test}}\}.
$$
The Neural ODE is then trained by replacing $\Omega$ with
$\Omega_{\mathrm{train}}$ in $\mathcal L_{\mathrm{ODE}}(\theta)$, while the test
MSE is computed as
$$
  \mathrm{MSE}_{\mathrm{test}}
  =
  \frac{1}{|\Omega_{\mathrm{test}}|}
  \sum_{(i,j)\in\Omega_{\mathrm{test}}}
  \left(y_i^{[j]}-\widehat y_i\right)^2.
$$
For these data, a natural forecasting split is to train on early times and test
on the last few time points, for example training through $t=132$ and testing
at $t=144,156,168$. With only 15 time points, an alternative is a rolling
leave-last-time-out protocol, such as training through $t=120$ and testing at
$t=132$, then training through $t=132$ and testing at $t=144$, and so on.

For AIC and BIC \citep{akaike1974newlook,schwarz1978estimating}, residuals were
modeled as independent Gaussian errors with
common variance $\sigma^2$. The maximum likelihood variance estimate is
$$
  \widehat\sigma^2 = \frac{\mathrm{SSE}}{n},
$$
which gives the maximized log-likelihood
$$
  \log \widehat L
  =
  -\frac{n}{2}
  \left[
    \log(2\pi)+1+\log(\widehat\sigma^2)
  \right].
$$
The criteria are
$$
    \mathrm{AIC}=2k-2\log\widehat L,
    \qquad
    \mathrm{BIC}=k\log(n)-2\log\widehat L
$$
Here $k=p+1$, where $p$ is the number of fitted dynamical parameters and the
additional parameter is the residual variance. Thus $k=3$ for the classical
logistic model, $k=1186$ for the Neural ODE, $k=2371$ for the Neural SDE,
$k=2212$ for the deterministic logistic PINN, and $k=2213$ for the stochastic
logistic PINN/SDE.
Initial conditions are not counted as fitted parameters because they are taken
directly from the first row of the data; for the single-curve Neural ODE this
means using $y_0=J^{-1}\sum_{j=1}^J y_0^{[j]}$.

The 90 percent interval coverage is computed for models that produce a
predictive distribution rather than only a fitted mean curve: the Neural SDE
and the stochastic logistic SDE. Let $B$ be the number of simulated SDE paths
($B=1000$ for the Neural SDE comparison and $B=5000$ for the stochastic
logistic SDE) and let
$Y_i^{(b)[j]}$ denote the simulated original-scale value for path $b$, time
$i$, and trial $j$. For interval level $\gamma=0.90$, set
$\alpha=(1-\gamma)/2=0.05$ and define the pointwise empirical interval
$$
  L_i^{[j]}=Q_{\alpha}\left(\{Y_i^{(b)[j]}\}_{b=1}^B\right),
  \qquad
  U_i^{[j]}=Q_{1-\alpha}\left(\{Y_i^{(b)[j]}\}_{b=1}^B\right),
$$
where $Q_q$ is the empirical $q$-quantile. The reported empirical coverage is
$$
  \widehat C_{0.90}
  =
  \frac{1}{n}
  \sum_{(i,j)\in\Omega}
  \mathbf 1\{L_i^{[j]}\le y_i^{[j]}\le U_i^{[j]}\}.
$$
Thus an observed data point is counted as covered if it lies between the 5
percent and 95 percent Monte Carlo quantiles for the same time and trial. For
the Neural SDE, the reported value $0.9333$ means that 42 of the 45 observed
responses fall inside their corresponding central 90 percent intervals. For
the stochastic logistic SDE, the reported value $0.8222$ means that 37 of the
45 observed responses fall inside their corresponding central 90 percent
intervals. The mean interval width is the average of $U_i^{[j]}-L_i^{[j]}$ over
the same observed entries. These are in-sample empirical calibration
diagnostics, not out-of-sample guarantees.